% ЗАКЛЮЧЕНИЕ — без номер. Ориентировъчен обем: 2–3 стр.
% По А.3 съдържа: обобщение на характеристиките и особеностите на решението,
% предимствата И НЕДОСТАТЪЦИТЕ му, резултатите от експериментите с тяхната
% инженерна трактовка, и предложения за по-нататъшна работа.
\section*{ЗАКЛЮЧЕНИЕ}
\label{sec:conclusion}
\addcontentsline{toc}{section}{ЗАКЛЮЧЕНИЕ}

В хода на дипломното проектиране е изградена работеща система, която описва графичен
потребителски интерфейс независимо от платформения инструментариум, и --- което е по-
същественото --- методика, с която това описание може да бъде \emph{проверено}, а не само
заявено. По задачите от Увода, в същия ред: направен е преглед на подходите и е показано
защо твърдението за еднакво изобразяване остава непроверено (Раздел~\ref{sec:literature});
анализирани са алтернативите по пет предварително изброени критерия и изборът е обоснован,
включително с преброяване в кодовата база вместо с преценка (Раздел~\ref{sec:alternatives});
проектирана е слоеста архитектура, в която платформено обусловената част е сведена до
таблично зададено съответствие (Раздел~\ref{sec:design}); тя е реализирана върху четири
native инструментариума и един слой без екран (Раздел~\ref{sec:implementation}); изградена
е методика за автоматизирано измерване на еквивалентността на изобразяването
(Раздел~\ref{sec:method}); и е проведена експериментална проверка върху \boardPages{}
интерфейса, от която е изведена систематизация на установените разминавания
(Раздел~\ref{sec:results}).

\textbf{Приноси.}
\begin{enumerate}[topsep=2pt,itemsep=2pt,leftmargin=*]
\item \textbf{Систематизация на нарушенията на независимостта} --- шест класа, всеки
      изведен от доведени до първопричина случаи, с обща констатация, която противоречи на
      очакването: абстракцията не протича на границата на програмния интерфейс, а на
      местата, където инструментариумът кодира неизказана уговорка. Това е приносът,
      който се обобщава извън конкретния артефакт.
\item \textbf{Методика за количествено измерване на еквивалентността} между две независими
      реализации на едно описание, включително правилата за оценяване, записани
      предварително, и механичните проверки за достоверност на самото измерване.
\item \textbf{Разграничение на средствата за откриване:} половината класове са откриваеми
      без екран, половината изискват сравнение на изображения --- от което следва, че
      измервателният апарат трябва да се изгражда успоредно със системата, а не след нея.
\item \textbf{Реализацията като инструмент} --- система върху пет платформени слоя, която
      прави измерването изобщо възможно, с публичен интерфейс, при който най-честият
      клас грешки при управление на паметта е грешка при компилация.
\item \textbf{Предварително регистрирани хипотези} за цената на реализацията, с изричните
      изисквания за валидност --- включително хипотеза, формулирана срещу очакването на
      автора.
\end{enumerate}

\textbf{Предимства на предложеното решение.}
Спрямо алтернативите от Раздел~\ref{sec:alternatives} решението печели по трите
критерия, които определят изхода на работата. По \textbf{К1 (проверимост)} --- послойното
изграждане, водено от тестове, е единствената от трите стратегии, при която съществува
общо описание, изпълнявано и от двете реализации, тоест контролно условие за сравнение. По
\textbf{К4 (безопасност по конструкция)} --- само преместваемият притежаващ манипулатор
превръща най-честия начин за създаване на цикъл в грешка при компилация, при измерена
ергономична цена, падаща върху малцинство от местата. По \textbf{К5 (цена за платформа)} ---
изборът на език дава пряк достъп до три от петте инструментариума без нито един ред
обвързващ код.

Отделно стои предимство, което не е било предвидено като цел. На настолната Apple
платформа еталонната реализация предлага \emph{само} пътя през мобилния инструментариум,
пренесен на настолната машина --- приложението изглежда като мобилно приложение в прозорец.
Разработената система носи \emph{и двата} слоя: същия мобилен път, който позволява прякото
сравнение в Раздел~\ref{sec:results}, и втори слой върху native настолния инструментариум,
чиито контроли са тези на операционната система. Изборът е на приложението и не изисква
промяна в описанието на интерфейса. Това е пряко следствие от архитектурата: щом
платформената част е сведена до таблица от действия, две различни таблици за една и съща
операционна система не са изключение, а обикновен случай.

Печели се и нещо друго, също непредвидено: побайтово съвпадение с еталона на част от
интерфейсите (§\ref{subsec:integrity_finding}) --- най-силната форма, която резултатът от
такова измерване изобщо може да приеме.

\textbf{Недостатъци и ограничения.}
\begin{enumerate}[topsep=2pt,itemsep=2pt,leftmargin=*]
\item \textbf{Четири от шестте регистрирани хипотези остават непроверени.} Измерванията за
      потребена памет и за време за реакция не са проведени; инструментът е проектиран, но
      не е изпълнен. Единствената проверена хипотеза се оказа \emph{твърде силна} по
      величина.
\item \textbf{Паритетът не е постигнат навсякъде.} На всяка платформа остават интерфейси с
      незначително разминаване, а на две --- и със съществено. Причините са отнесени към
      класове в §\ref{subsec:taxonomy}, но не са отстранени.
\item \textbf{Един платформен слой е реализиран частично по замисъл} --- изобразяват се
      само част от контролите, останалите използват огледалото без екран.
\item \textbf{Цената на модела на собствеността} е реална: достъпът до задържана контрола
      изисква проверка за живост, а съвместното притежаване --- изрично действие. Това е
      съзнателна размяна на удобство срещу проверимост, не безплатно предимство.
\item \textbf{Инструментът не може да наблюдава освобождаване на памет} при напускане на
      интерфейс, тъй като всяко измерване се извършва в отделен процес
      (§\ref{subsec:capture}) --- ограничение на конструкцията, не пропуск в изпълнението.
\item \textbf{Резултатът е обвързан с една версия на еталона} и не се пренася автоматично
      към следваща.
\item \textbf{Един критерий за достоверност остаря заради успеха на системата}
      (§\ref{subsec:integrity_finding}) --- проверката обявява побайтовото съвпадение за
      дефект, което е било вярно, докато то е било недостижимо.
\end{enumerate}

\textbf{Насоки за по-нататъшна работа.}
\begin{enumerate}[topsep=2pt,itemsep=2pt,leftmargin=*]
\item \textbf{Провеждане на измерванията по H2a и H2b} --- пряко следствие от ограничение~1.
      Инструментът е проектиран; остава изпълнението и отчитането, включително на
      отрицателния резултат, ако H2b-3 се потвърди.
\item \textbf{Стесняване на критерия за достоверност} --- следствие от ограничение~7:
      признакът трябва да стане „междуколонно съвпадение \emph{и} съвпадение между темите“,
      иначе броят лъжливи тревоги ще расте с качеството на реализацията.
\item \textbf{Довършване на частичния платформен слой} (ограничение~3) и повторно
      измерване --- това е и проверка дали систематизацията от §\ref{subsec:taxonomy}
      предсказва правилно кои класове ще се появят на нов инструментариум.
\item \textbf{Платформен слой за Linux} --- най-информативното възможно продължение,
      защото еталонната реализация \emph{няма} такъв. Причината не е техническа
      невъзможност, а следствие от собствената ѝ архитектура: тя не рисува сама, а
      преобразува към native инструментариум, и затова изисква на всяка платформа да
      съществува \emph{един} инструментариум, чиито контроли са очевидният избор ---
      какъвто на Linux няма (GTK и Qt съжителстват, а работната среда не предопределя
      кой от двата е „системният“). Към това се добавя, че поддръжката на платформа за
      нея означава и разпространение на средата за изпълнение за нея.

      Разработената система няма нито едното ограничение: тя не носи среда за изпълнение,
      а слоят се пише срещу избран инструментариум. Основното предизвикателство се измества
      другаде --- при липса на еталон на тази платформа \emph{отпада оракулът}, а с него и
      методиката от Раздел~\ref{sec:method}. Проверката там може да бъде само вътрешна:
      съгласуваност между двата входни пътя и съответствие с изчисленото от преносимия слой
      разположение, но не и съвпадение с външен образец. Това е и границата на самия метод,
      която новата платформа би направила видима.
\item \textbf{Разширяване на систематизацията с шести платформен слой} от друго семейство ---
      всяка нова среда е \emph{повторение} на експеримента и проверява дали класовете са
      свойство на задачата или на конкретните четири инструментариума.
\item \textbf{Автоматизирано следене на разминаванията във времето} --- запазване не само
      на текущото състояние, а и на заснеманията при всяка кампания, което би направило
      възможна и кривата от §\ref{subsec:agreement}, и съпоставките ПРЕДИ/СЛЕД, които
      сега липсват (ограничение по §\ref{subsec:visual}).
\end{enumerate}
